\documentclass[11pt,a4paper]{article}

\usepackage[a4paper,margin=2.7cm]{geometry}
\usepackage{amsmath,amssymb,amsthm,mathtools,mathrsfs}
\usepackage{microtype}
\usepackage[hidelinks]{hyperref}
\usepackage{enumitem}

\newtheorem{theorem}{Theorem}
\newtheorem{lemma}[theorem]{Lemma}
\newtheorem{corollary}[theorem]{Corollary}
\newtheorem{proposition}[theorem]{Proposition}
\newtheorem{conjecture}[theorem]{Conjecture}
\theoremstyle{definition}
\newtheorem{definition}[theorem]{Definition}
\newtheorem{problem}[theorem]{Problem}
\theoremstyle{remark}
\newtheorem{remark}[theorem]{Remark}

\newcommand{\E}{\mathbb E}
\newcommand{\I}{\mathcal I}
\newcommand{\T}{\mathcal T}
\newcommand{\Sset}{\mathscr S}
\newcommand{\1}{\mathbf 1}
\newcommand{\dd}{\mathrm d}

\title{Finite--Zero Temperature Bellman Duality\\
for Extremal Independent Sets on Trees}
\author{Research note}
\date{2 October 2026}

\begin{document}
\maketitle

\begin{abstract}
For a finite tree $T$ and hard-core activity $\lambda>0$, let
\[
  \mu_\lambda(T)=\frac{\lambda Z_T'(\lambda)}{Z_T(\lambda)},
  \qquad
  a_*(\lambda)=\inf_T \frac{\mu_\lambda(T)}{\alpha(T)}.
\]
We give an exact local representation of both the finite-temperature observable
$\mu_\lambda$ and the zero-temperature observable $\alpha$ on the same rooted tree.
Each vertex carries a three-component state $(x,z,g)$ satisfying a closed recursion,
and
\[
  \mu_\lambda(T)=\sum_{v\in T} z_v,
  \qquad
  \alpha(T)=\sum_{v\in T} g_v,
  \qquad g_v\in\{0,1\}.
\]
This yields an exact Bellman dual characterization of $a_*(\lambda)$: it is the largest
$r$ for which there exists a nonpositive potential on the reachable state space whose
one-step Bellman slack is nonnegative for every rooted-tree composition.  The duality
is exact in both directions and immediately supplies a rigidity principle: near-extremal
trees must concentrate on the contact set of any sharp dual potential.

At $\lambda=1$ the child multiset enters the parent state only through three aggregate
statistics, making the problem suitable for a rigorous interval-certified Bellman search.
The concrete target is a certificate at $r=0.4733$.  Combined with the existing explicit
upper witness $a_*(1)<0.47330186$, such a certificate would confine the constant to an
interval of width $1.86\times10^{-6}$.
\end{abstract}

\section{The extremal constant}

For a finite nonempty tree $T$, write
\[
  Z_T(\lambda)=\sum_{I\in\I(T)}\lambda^{|I|},
  \qquad
  \mu_\lambda(T)=\lambda\frac{Z_T'(\lambda)}{Z_T(\lambda)},
\]
and let $\alpha(T)$ be the independence number.  Define
\begin{equation}
  a_*(\lambda)
  :=\inf_{T\text{ finite nonempty tree}}
  \frac{\mu_\lambda(T)}{\alpha(T)}.
  \label{eq:astar}
\end{equation}
The purpose of this note is to replace the optimization over all finite trees by an exact
local dynamic program admitting a dual certificate.

\section{A common local state for two temperatures}

Root a finite tree $T$.  For a vertex $v$, let $T_v$ denote the rooted descendant subtree,
and let $C(v)$ be the children of $v$.

For the hard-core model, write $Z_{0,v}$ and $Z_{1,v}$ for the partition functions of
$T_v$ with $v$ forced vacant and occupied, respectively, and define the cavity ratio
\begin{equation}
  x_v:=\frac{Z_{1,v}}{Z_{0,v}}.
  \label{eq:xdef}
\end{equation}
Let $D:=\lambda\partial_\lambda$ and define the local response charge
\begin{equation}
  z_v:=D\log(1+x_v).
  \label{eq:zdef}
\end{equation}

For the zero-temperature problem, let $\alpha_{0,v}$ be the largest independent-set size
in $T_v$ subject to $v$ being excluded, and put
\begin{equation}
  g_v:=\alpha(T_v)-\alpha_{0,v}.
  \label{eq:gdef}
\end{equation}

\begin{theorem}[Exact finite/zero-temperature recursion]
\label{thm:local-recursion}
For every rooted finite tree and every $\lambda>0$,
\begin{align}
  x_v
  &=\lambda\prod_{u\in C(v)}(1+x_u)^{-1},
  \label{eq:xrec}\\
  z_v
  &=\frac{x_v}{1+x_v}
    \left(1-\sum_{u\in C(v)}z_u\right),
  \label{eq:zrec}\\
  g_v
  &=\1\!\left\{g_u=0\text{ for every }u\in C(v)\right\}.
  \label{eq:grec}
\end{align}
In particular $g_v\in\{0,1\}$ for every vertex.
\end{theorem}

\begin{proof}
The usual root decomposition gives
\[
  Z_{0,v}=\prod_{u\in C(v)}(Z_{0,u}+Z_{1,u}),
  \qquad
  Z_{1,v}=\lambda\prod_{u\in C(v)}Z_{0,u},
\]
which proves \eqref{eq:xrec}.  Taking $D\log$ gives
\[
  D\log x_v
  =1-\sum_{u\in C(v)}D\log(1+x_u)
  =1-\sum_{u\in C(v)}z_u.
\]
Since
\[
  D\log(1+x_v)=\frac{x_v}{1+x_v}D\log x_v,
\]
we obtain \eqref{eq:zrec}.

For the independence number,
\[
  \alpha_{0,v}=\sum_{u\in C(v)}\alpha(T_u),
\]
whereas forcing $v$ into the independent set gives
\[
  \alpha_{1,v}=1+\sum_{u\in C(v)}\alpha_{0,u}.
\]
Hence
\[
  \alpha_{1,v}-\alpha_{0,v}
  =1-\sum_{u\in C(v)}g_u.
\]
By induction from a leaf, each $g_u$ is $0$ or $1$, and therefore
\[
  g_v
  =\max\left\{0,1-\sum_{u\in C(v)}g_u\right\}
  =\1\!\left\{\sum_{u\in C(v)}g_u=0\right\}.
\]
\end{proof}

The remarkable point is that the two global observables themselves telescope.

\begin{theorem}[Exact local-charge decomposition]
\label{thm:charges}
For every rooted finite tree $T$,
\begin{equation}
  \boxed{
  \mu_\lambda(T)=\sum_{v\in T}z_v,
  \qquad
  \alpha(T)=\sum_{v\in T}g_v.}
  \label{eq:charges}
\end{equation}
Moreover
\[
  A_g:=\{v:g_v=1\}
\]
is a maximum independent set of $T$.
\end{theorem}

\begin{proof}
Since
\[
  Z_{T_v}
  =Z_{0,v}(1+x_v)
  =\left(\prod_{u\in C(v)}Z_{T_u}\right)(1+x_v),
\]
we have
\[
  \mu_\lambda(T_v)
  =\sum_{u\in C(v)}\mu_\lambda(T_u)+z_v.
\]
Induction proves the first identity in \eqref{eq:charges}.

Likewise, by definition of $g_v$,
\[
  \alpha(T_v)=\alpha_{0,v}+g_v
  =\sum_{u\in C(v)}\alpha(T_u)+g_v,
\]
which proves the second identity.  If $g_v=1$, then every child has $g_u=0$ by
\eqref{eq:grec}, so $A_g$ is independent.  Its cardinality is $\sum_v g_v=\alpha(T)$,
hence it is maximum.
\end{proof}

\begin{remark}
The set $A_g$ is a canonical maximum independent set attached to the chosen rooting.
Thus the finite-temperature hard-core field and a zero-temperature maximum independent set
are encoded on exactly the same recursion tree.
\end{remark}

\section{Exact Bellman duality}

Let $\Sset_\lambda$ be the set of all states
\[
  s(T)=(x(T),z(T),g(T))
\]
that occur at the root of a finite rooted tree at activity $\lambda$.  If
$s_1,\dots,s_d\in\Sset_\lambda$ are child states, write
\[
  F_\lambda(s_1,\dots,s_d)=s
\]
for the parent state determined by \eqref{eq:xrec}--\eqref{eq:grec}.  The case $d=0$
is the leaf state
\begin{equation}
  s_\bullet
  =\left(\lambda,\frac{\lambda}{1+\lambda},1\right).
  \label{eq:leaf}
\end{equation}

For $r\in\mathbb R$, define the local cost
\begin{equation}
  c_r(x,z,g):=z-rg.
  \label{eq:local-cost}
\end{equation}

\begin{definition}[Bellman certificate]
A function $\Phi:\Sset_\lambda\to\mathbb R$ is an $r$-certificate if
\begin{equation}
  \Phi(s)\le0
  \qquad(s\in\Sset_\lambda)
  \label{eq:nonpositive}
\end{equation}
and for every finite child list and its parent state
$s=F_\lambda(s_1,\dots,s_d)$,
\begin{equation}
  c_r(s)+\Phi(s)-\sum_{i=1}^d\Phi(s_i)\ge0.
  \label{eq:bellman-ineq}
\end{equation}
\end{definition}

\begin{theorem}[Bellman lower-bound theorem]
\label{thm:bellman-lb}
If an $r$-certificate exists, then
\[
  a_*(\lambda)\ge r.
\]
More precisely, every finite tree satisfies
\[
  \mu_\lambda(T)-r\alpha(T)\ge-\Phi(s(T))\ge0.
\]
\end{theorem}

\begin{proof}
Apply \eqref{eq:bellman-ineq} at every vertex.  Every nonroot potential appears once
with a plus sign at that vertex and once with a minus sign as a child of its parent.
Therefore
\[
  \sum_{v\in T}\left(z_v-rg_v\right)+\Phi(s(T))\ge0.
\]
Theorem~\ref{thm:charges} turns the sum into
$\mu_\lambda(T)-r\alpha(T)$.  Finally use $\Phi\le0$.
\end{proof}

The converse is equally important: this is not merely a sufficient device.

\begin{theorem}[Exact Bellman duality]
\label{thm:bellman-duality}
For every $\lambda>0$,
\begin{equation}
  \boxed{
  a_*(\lambda)
  =\max\Bigl\{r\in\mathbb R:
  \text{an $r$-certificate on $\Sset_\lambda$ exists}\Bigr\}.}
  \label{eq:dual-characterization}
\end{equation}
\end{theorem}

\begin{proof}
Theorem~\ref{thm:bellman-lb} proves that the right-hand side is at most $a_*(\lambda)$.

Conversely fix $r\le a_*(\lambda)$.  For a reachable state $s$, define
\begin{equation}
  V_r(s)
  :=\inf\left\{
      \mu_\lambda(T)-r\alpha(T):
      T\text{ finite rooted tree with }s(T)=s
    \right\}.
  \label{eq:value}
\end{equation}
Every term inside the infimum is nonnegative, so $V_r(s)\ge0$, and the set over which
we take the infimum is nonempty by reachability.  Put
\[
  \Phi_r(s):=-V_r(s)\le0.
\]

Take any transition $s=F_\lambda(s_1,\dots,s_d)$.  For $\varepsilon>0$, choose rooted
trees $T_i$ with state $s_i$ and
\[
  \mu_\lambda(T_i)-r\alpha(T_i)
  \le V_r(s_i)+\varepsilon/d
\]
(with the obvious interpretation when $d=0$).  Attach their roots to a new parent.
By Theorem~\ref{thm:charges}, the resulting tree has state $s$ and deficit
\[
  c_r(s)+\sum_{i=1}^d
  \bigl(\mu_\lambda(T_i)-r\alpha(T_i)\bigr).
\]
Hence
\[
  V_r(s)
  \le c_r(s)+\sum_{i=1}^d V_r(s_i)+\varepsilon.
\]
Letting $\varepsilon\downarrow0$ and substituting $\Phi_r=-V_r$ gives
\eqref{eq:bellman-ineq}.  Thus an $r$-certificate exists for every $r\le a_*(\lambda)$,
including $r=a_*(\lambda)$ itself.
\end{proof}

\begin{remark}[Why this is useful]
The proof of Theorem~\ref{thm:bellman-duality} constructs an abstract optimal potential,
but not an explicit one.  The research problem is therefore sharply separated into two parts:
\begin{enumerate}[label=(\roman*)]
  \item the extremal constant is exactly a Bellman dual optimum;
  \item proving a numerical lower bound reduces to exhibiting one explicit potential.
\end{enumerate}
There is no loss from the dual formulation.
\end{remark}

\section{Slack identity and rigidity}

Given an $r$-certificate $\Phi$, define the one-step slack
\begin{equation}
  \Delta_\Phi(s;s_1,\dots,s_d)
  :=c_r(s)+\Phi(s)-\sum_{i=1}^d\Phi(s_i)\ge0.
  \label{eq:slack}
\end{equation}

\begin{proposition}[Exact slack decomposition]
\label{prop:slack}
Every rooted finite tree satisfies
\begin{equation}
  \mu_\lambda(T)-r\alpha(T)
  =\sum_{v\in T}\Delta_\Phi(v)-\Phi(s(T)),
  \label{eq:slack-decomp}
\end{equation}
where $\Delta_\Phi(v)$ denotes the slack of the transition at $v$.
Both terms on the right are nonnegative.
\end{proposition}

\begin{proof}
This is the same telescoping computation as in
Theorem~\ref{thm:bellman-lb}, now retaining the exact slack at each vertex.
\end{proof}

\begin{corollary}[Contact-set rigidity]
\label{cor:rigidity}
Assume $r=a_*(\lambda)$ and let $T_n$ be a sequence with
\[
  \frac{\mu_\lambda(T_n)}{\alpha(T_n)}\longrightarrow r,
  \qquad
  \alpha(T_n)\longrightarrow\infty.
\]
If an optimal certificate $\Phi$ additionally satisfies
\[
  -\Phi(s(T_n))=o(\alpha(T_n)),
\]
then
\[
  \sum_{v\in T_n}\Delta_\Phi(v)=o(\alpha(T_n)).
\]
Consequently, any family of local transitions on which
$\Delta_\Phi\ge\eta>0$ can occur only $o(\alpha(T_n))$ times.
\end{corollary}

This turns an explicit sharp dual potential into an extremal classification theorem: asymptotically
optimal trees must live on the contact set $\Delta_\Phi=0$.

\section{Specialization to $\lambda=1$}

At unit activity the leaf state is
\begin{equation}
  s_\bullet=(1,1/2,1),
\end{equation}
and the transition becomes
\begin{align}
  x&=\prod_{i=1}^d(1+x_i)^{-1},
  \label{eq:unit-x}\\
  z&=\frac{x}{1+x}\left(1-\sum_{i=1}^d z_i\right),
  \label{eq:unit-z}\\
  g&=\1\!\left\{\sum_{i=1}^d g_i=0\right\}.
  \label{eq:unit-g}
\end{align}

There is an important compression.  Put
\begin{equation}
  L:=\sum_i\log(1+x_i),
  \qquad
  Z:=\sum_i z_i,
  \qquad
  K:=\sum_i g_i.
  \label{eq:aggregates}
\end{equation}
Then
\begin{equation}
  \boxed{
  x=e^{-L},
  \qquad
  z=\frac{x}{1+x}(1-Z),
  \qquad
  g=\1\{K=0\}.}
  \label{eq:aggregate-transition}
\end{equation}
Thus an unbounded child multiset influences the parent state only through two real additive
statistics and the zero/nonzero status of one integer statistic.

\subsection{A coercive potential ansatz}

A useful explicit family is
\begin{equation}
  \Phi(x,z,g)
  =-A\log(1+x)-Bg-H_g(x,z),
  \qquad A,B\ge0,\quad H_g\ge0.
  \label{eq:ansatz}
\end{equation}
It is automatically nonpositive.  If $K=\sum_i g_i$, then the local slack at $\lambda=1$
is exactly
\begin{align}
  \Delta_\Phi
  ={}&z-rg
  -A\log\bigl(x(1+x)\bigr)
  +B(K-g)
  \nonumber\\
  &\quad
  -H_g(x,z)+\sum_i H_{g_i}(x_i,z_i).
  \label{eq:ansatz-slack}
\end{align}
The first line is explicit and contains two useful coercive terms:
$B K$ penalizes many zero-temperature active children, while the logarithmic term records the
entire multiplicative hard-core recursion exactly.  The remaining task is to construct two
nonnegative correction functions $H_0,H_1$.

Equation \eqref{eq:ansatz-slack} is the proposed route to a machine-checkable proof: approximate
$H_0,H_1$ by piecewise-rational or piecewise-polynomial functions, prove the Bellman inequality
by interval arithmetic on boxes, and handle the unbounded child count through the additive
statistics in \eqref{eq:aggregates} rather than by an arity cutoff.

\section{The primal laboratory: layered trees}

The known strong upper witnesses are recursively symmetric, so it is useful to record the
one-dimensional primal dynamics.  Start from
\[
  x_0=1,
  \quad z_0=\frac12,
  \quad g_0=1,
  \quad M_0=\frac12,
  \quad A_0=1,
  \quad n_0=1.
\]
If level $j+1$ consists of a new root attached to $b_{j+1}\ge1$ identical copies of the
previous rooted tree, then
\begin{align}
  x_{j+1}&=(1+x_j)^{-b_{j+1}},
  \label{eq:layer-x}\\
  z_{j+1}&=\frac{x_{j+1}}{1+x_{j+1}}
              \left(1-b_{j+1}z_j\right),
  \label{eq:layer-z}\\
  g_{j+1}&=1-g_j,
  \label{eq:layer-g}\\
  M_{j+1}&=b_{j+1}M_j+z_{j+1},
  \label{eq:layer-mu}\\
  A_{j+1}&=b_{j+1}A_j+g_{j+1},
  \label{eq:layer-alpha}\\
  n_{j+1}&=1+b_{j+1}n_j.
  \label{eq:layer-n}
\end{align}
Here $M_j=\mu(T_j)$ and $A_j=\alpha(T_j)$ exactly.

An existing exact witness uses the bottom-up branching sequence
\begin{equation}
  (2,1,3,9,6,9,8,5)
  \label{eq:known-sequence}
\end{equation}
and gives
\begin{equation}
  a_*(1)<0.47330186.
  \label{eq:known-upper}
\end{equation}
The layered dynamics should be viewed as a primal search engine; the Bellman potential is its
matching dual certificate.

\section{Compatibility with the dilute golden-ratio regime}

For a star with $d$ leaves, all leaves have
\[
  x=\lambda,
  \qquad
  z=q:=\frac{\lambda}{1+\lambda},
  \qquad
  g=1.
\]
The center therefore has
\[
  x_c=\frac{\lambda}{(1+\lambda)^d},
  \qquad
  z_c=\frac{x_c}{1+x_c}(1-dq),
  \qquad
  g_c=0,
\]
and hence
\begin{equation}
  \frac{\mu_\lambda(K_{1,d})}{\alpha(K_{1,d})}
  =q+\frac{z_c}{d}.
  \label{eq:star-exact}
\end{equation}
Taking $d\sim s/\lambda$ gives
\[
  \frac{\mu_\lambda}{\alpha}
  =q-\left(1-\frac1s\right)e^{-s}\lambda^2+o(\lambda^2).
\]
The maximizing $s>1$ satisfies
\[
  s^2-s-1=0,
\]
so $s=\varphi=(1+\sqrt5)/2$ and the coefficient is
$e^{-\varphi}/\varphi^2$.  Thus the golden-ratio asymptotic appears directly as a stationary
point of the same local dynamics used in the Bellman formulation.

\section{Relation to the entropy decomposition at $\lambda=1$}

The canonical set $A_g=\{v:g_v=1\}$ from Theorem~\ref{thm:charges} is a maximum independent
set, so the entropy--partition decomposition can be applied with this specific $A_g$.
The existing scalar argument yields
\begin{equation}
  c
  =\frac{\log(34/5)}{6\log2}
  =0.4609224577271628\ldots
  \le a_*(1).
  \label{eq:entropy-lb}
\end{equation}
Its exact nonnegative decomposition separates four sources of slack: a degree term, an overlap
term, a coordinate-comparison term, and a Kullback--Leibler term.  The scalar degree inequality
is tight only at degree four, but that does not force the other slacks to vanish.

The Bellman state $(x,z,g)$ retains information discarded by the scalar degree relaxation.
Accordingly, the entropy lower bound should be regarded as a coarse feasible dual point, while
Theorem~\ref{thm:bellman-duality} identifies the exact dual optimum.

\section{The concrete $\lambda=1$ target}

The present rigorous interval is
\begin{equation}
  0.4609224577271628\ldots
  \le a_*(1)
  <0.47330186.
  \label{eq:current-interval}
\end{equation}
The proposed first target is the rational value
\begin{equation}
  r_0:=\frac{4733}{10000}=0.4733.
  \label{eq:r0}
\end{equation}

\begin{problem}[Finite certificate target]
Construct explicit nonnegative functions $H_0,H_1$ and constants $A,B\ge0$ such that the
potential \eqref{eq:ansatz} satisfies
\begin{equation}
  \Delta_\Phi\ge0
  \label{eq:target-ineq}
\end{equation}
for every reachable unit-activity transition, with $r=r_0$.
\end{problem}

By Theorem~\ref{thm:bellman-lb}, such a certificate would imply
\begin{equation}
  \boxed{
  0.4733\le a_*(1)<0.47330186,}
  \label{eq:bomb-interval}
\end{equation}
an interval of width $1.86\times10^{-6}$.

This is a qualitatively different task from another global entropy estimate.  It is a local
verification problem on the exact finite/zero-temperature state recursion.

\section{A certification architecture}

A rigorous computer-assisted proof can be organized as follows.

\begin{enumerate}[label=\textbf{C\arabic*.},leftmargin=2.2em]
\item \textbf{Reachable-state enclosure.}
Generate certified outer boxes for the two reachable sectors $g=0$ and $g=1$ under
\eqref{eq:unit-x}--\eqref{eq:unit-g}.  The enclosure is recursive and uses directed rational
or interval arithmetic only.

\item \textbf{Potential representation.}
Represent $H_0,H_1$ by a piecewise-rational, piecewise-affine, or Bernstein-polynomial
function with rational coefficients and explicit interval remainder bounds.

\item \textbf{Aggregate child reduction.}
Use \eqref{eq:aggregates} to eliminate the raw child count from the parent equations.  The
remaining optimization of $\sum_i H_{g_i}(x_i,z_i)$ is an additive moment problem under
fixed $(L,Z,K)$.

\item \textbf{Tail coercivity.}
Choose the explicit $A$ and $B$ terms in \eqref{eq:ansatz} so that configurations with large
$L$ or large $K$ have a strictly positive analytic margin.  Only a compact residual region is
then passed to interval subdivision.

\item \textbf{Exact checker.}
The final verifier receives only rational box endpoints and rational potential coefficients.
Every transcendental quantity is enclosed by monotone rational bounds; the checker verifies
\eqref{eq:target-ineq} box by box and prints the minimum certified slack.
\end{enumerate}

The desired end product is therefore not a stochastic search result but a short deterministic
certificate whose mathematical statement is precisely Theorem~\ref{thm:bellman-lb}.

\section{What is proved here and what remains}

The following statements are proved in this note without computational assumptions:
\begin{enumerate}[label=(\roman*)]
  \item the common state recursion \eqref{eq:xrec}--\eqref{eq:grec};
  \item the exact charge identities \eqref{eq:charges};
  \item the canonical maximum independent set $A_g$;
  \item the Bellman lower-bound theorem;
  \item the converse and exact dual characterization \eqref{eq:dual-characterization};
  \item the exact slack/rigidity identity \eqref{eq:slack-decomp};
  \item the aggregate compression \eqref{eq:aggregate-transition} at $\lambda=1$.
\end{enumerate}

The only missing step toward the interval \eqref{eq:bomb-interval} is an explicit verified
potential satisfying \eqref{eq:target-ineq} at $r=0.4733$.  That is now a sharply formulated
finite-temperature/zero-temperature dual-certificate problem rather than an unconstrained
search over trees.

\end{document}
